We study whether multilingual sycophantic capitulation rises as language resource level falls. The motivating hypothesis is straightforward: if factual representations are weaker in lower-resource languages, challenge-induced agreement pressure may more easily override them. To test this claim, we build a seven-language evaluation that combines a multilingual extension of \bullshitbench with a 20-item \mkqa subset across English, French, Arabic, Hindi, Swahili, Yoruba, and Tagalog. We evaluate \gptmini and \qwen, then run a mechanistic follow-up on Qwen with hidden-state contrast analysis and activation steering.

The main result is negative in an important way. Lower-resource languages consistently reduce factual robustness, but they do not produce higher false-premise capitulation in our challenge setting. On \mkqa, the gap between high- and low-resource tiers is 20.0 points for \gptmini and 22.5 points for \qwen in final accuracy. On \bullshitbench, however, high-resource languages show equal or higher capitulation than low-resource languages, and high-vs-low Fisher tests are not significant for either model. The strongest intervention result comes from the local mechanistic analysis: activation steering on Qwen raises final rejection from 0.693 to 0.947 and reduces measured capitulation from 0.0129 to 0.0000.

These findings revise the original hypothesis. The multilingual risk exposed here is a split failure profile: lower-resource languages weaken factual knowledge, while false-premise prompts can elicit more conservative rejection rather than more agreement. This makes benchmark design, not just language coverage, central to future multilingual sycophancy evaluation.
